% !TEX root = ../nc-cubic.tex
\subsection{三重積の核による二種類の標的の復元}
\label{nft:sec:triple-kernel}
\begin{theorem}\label{nft:thm:triple-kernel}
$\C$ 上で次の三種類の dg 圏の合併を考える。滑らかな可換三次五様体の
$\Ku$、滑らかな次数一 del Pezzo 三様体の $\Ku$、および
系\ref{nft:cor:allsigns-five}の滑らかな偶符号被覆の残余圏である。
$\rho=S^{-1}[3]$、$R_j=\operatorname{Hom}(\mathrm{id},\rho^j)$
（対角双加群間の次数零の射）、$V=R_1$ と置く。
このクラスでは $V$ の積は可換で、三重積の核
\[
 I_3=\ker\big(\operatorname{Sym}^3V\longrightarrow R_3\big)
\]
の線型同型類 $(V,I_3)$ が dg Morita 同値類を決定する。
可換三次五様体で実現される条件は三重積の階数が $35$ であることで、
残る場合の階数は $34$ である。二重積だけではこのクラスを分類できない。
\end{theorem}
\begin{proof}
偶符号被覆には系\ref{nft:cor:allsigns-five}と
定理\ref{nft:thm:sign-sextic-kernel}の比較双加群を適用する。
可換三次側の一次生成環は七変数の滑らかな三次式の Jacobian 環である
（命題\ref{nft:prop:jacobian-necessary}）。もう一方の標的は平方完成と
重み二の座標の平行移動により
\[
 Z_H=(s^2+H=0),\qquad H=w^3+wB_4(v)+B_6(v)
\]
の形を取る。$S_0=\C[v_1,v_2,v_3,w]$、$\deg v_i=1,\ \deg w=2$ と置き、$J=(H_w,H_{v_1},H_{v_2},H_{v_3})\subset S_0$ とする。
$\rho=(2)$ に対する対角核の固定点計算は
系\ref{nft:cor:sign-full-torelli}直後と同一である。
非恒等 sector は固定座標 $s$ または $w$ での制限 $s^2,w^3$ と
法線の重みしか用いず、次数一、二に現れない。恒等 sector の積は
多項式の積である \cite[Theorem~5.39 の証明]{NFBFKKernels}。
従って $V=(S_0)_2$ で、一次生成環は $(S_0/J)_{\rm even}$ となる。
実際 $\operatorname{Sym}^j V\to (S_0)_{2j}$ は全射である。
この偶数部分の同定は \cite[Theorem~4.3]{SignLRZVeronese} とも一致する。

両側の一次生成環は socle 次数七の完全対を持つ。三次側は七本の二次の
正則列、六次側は孤立性による重み $4,5,5,5$ の正則列を用いる。
後者の socle は重み $14$ で、完全対は偶数部分に制限しても完全である。
よって三重積の核から二重積の核は
\[
 I_2=\{f\in\operatorname{Sym}^2V: Vf\subset I_3\}
\]
と復元できる。逆包含は次数二の類が全 $V$ に消されれば socle に入ることによる。
三次側の次数三の次元は $(1+t)^7$ の三次係数 $35$。
六次側の重み六の単項式は $50$ 個で、$J_6$ は $H_w(S_0)_2$ と
$\sum_i H_{v_i}(S_0)_1$ による $7+9=16$ 本で張られる。
正則列の最初の Koszul 関係の重みは九なのでこれらは独立で、差は $34$ となる。
三次側では $I_2$ が全偏微分の空間であり、
命題\ref{gcd:prop:integrability-symmetrizer}により積分して滑らかな三次式を
射影同型を除き復元できる。

六次側での座標の復元には $R_4$ も非恒等 sector の元も必要としない。
$D=\operatorname{Proj}(\operatorname{Sym}V/(I_2))$ を埋込とともに取る。
$(S_0)_{\rm even}$ の表示の核は対称行列 $(v_iv_j)$ の二次小行列式で、
追加の二次関係は $q=H_w=3w^2+B_4$ のみである。従って
$D$ は重み付き射影空間での $q=0$ の二次 Veronese 表示である。
これは頂点を避け、$\pi:D\to\mathbb P^2$ は $\mathcal O(2)$ の全空間で
モニック二次式を取った有限平坦二重被覆である。従って
\[
 \pi_*\mathcal O_D=\mathcal O\oplus\mathcal O(-2),\qquad
 \omega_D=\pi^*\mathcal O(-1),\qquad
 \mathcal O_D(1)_{\rm emb}=\pi^*\mathcal O(2).
\]
中央の式は全空間の標準束 $\pi^*\mathcal O(-5)$ と方程式の重み四の
随伴公式による。$B_4=0$ の非被約二重平面を含めて成り立つ。
制限 $V\to H^0(\mathcal O_D(1)_{\rm emb})$ は同型であり、
$H^0(\mathcal O_D)=\C$ より $\omega_D^{-2}\simeq\mathcal O_D(1)_{\rm emb}$
の選択はスカラー倍だけである。そこで積
\[
 \operatorname{Sym}^2H^0(D,\omega_D^{-1})\longrightarrow H^0(D,\omega_D^{-2})
\]
の像を取ると $V$ の中の六次元部分空間 $\langle v_iv_j\rangle$ を内在的に得る。
三次側では $I_2$ が正則列なので対応する $D$ は空であり、
この空性だけでも二種の標的を区別できる。

$H^0(D,\omega_D^{-1})$ の基底と上記の像の補方向 $w$ を選ぶことで、
$V=(S_0)_2$ の表示を得る。選択の差は $v\mapsto Mv$、
$w\mapsto aw+Q_2(v)$ とスカラーの平方根で尽くされる。
既知となった $\operatorname{Sym}^jV\to(S_0)_{2j}$ による
$I_2,I_3$ の像は $J_4,J_6$ である。残る生成重み五は
\[
 J_5=\{h\in(S_0)_5:v_ih\in J_6\ (i=1,2,3)\}
\]
と復元する。実際 $S_0/J$ の完全対により全 $v_i$ の同時零化空間は
$S_0/(J,v_1,v_2,v_3)=\C[w]/(w^2)$ の socle 重み $14$ に関する双対で、
重み五の成分は消える。従って $J=(J_4,J_5)$ が決まる。
重み六の未知式 $F$ について $F_w\in J_4$, $F_{v_i}\in J_5$ を解き、
偏微分が $J/(S_0)_+J$ を張る解を取る。元の $H$ が存在を保証する。
系\ref{nft:cor:sign-sextic}の鉛筆と重み付きベクトル場による証明は、
正規化した二つの六次式の勾配イデアルの一致だけを使うので、
得られる滑らかな方程式は重み付き座標変換を除き一意である。

復元した標的の座標同型のグラフを既述の比較双加群と合成すると実同値を得る。
逆に dg Morita 同値は Serre 関手とシフト、従って次数一からの三重合成を運ぶ。
最後に系\ref{nft:cor:hesse-nonrealizable}の切片はすべて $q=3w^2$ を持つ。
その二重積は恒等 sector への全射で、核は上記の同じ小行列式と $w^2$ である。
ところが $\kappa^3$ の異なる二点の圏は同値でない。これで次数の最適性を得る。
\end{proof}
以上は指定クラスに所属する圏に対する判定である。未知の滑らかな dg 圏から
同じ三重積が得られる場合に、その圏と復元多様体との比較核を与える問題は残る。

\begin{corollary}[指定二因子族と零係数偶符号族の識別]
\label{nft:cor:triple-family}
定理\ref{nft:thm:triple-kernel}の合併の中で、固定した二つの Sklyanin 因子の
一中心族（系\ref{tes:cor:hyperplane-recognition}の仮定）への所属は、
三重積の階数が $35$ であり、積分した三次式が同系の二平面条件を満たすことと
同値である。一方、滑らかな零係数偶符号族への所属は、階数が $34$ で、
復元して規格化した $H=w^3+wB_4+B_6$ の $B_4,B_6$ が同じ三つの
一次座標の平方の多項式で書けることと同値である。所属は dg Morita 同値を除いて取る。
\end{corollary}
\begin{proof}
階数 $35$ の場合、定理の比較双加群により二平面条件の判定に必要な
可換三次との同値が得られ、系\ref{tes:cor:hyperplane-recognition}を適用できる。
階数 $34$ の場合、同定理で得た $Z_H$ に系\ref{nft:cor:sign-target-locus}を適用する。
どちらも標的のグラフと各族の比較核を合成して十分性を得る。
\end{proof}

\begin{corollary}[二重積が記憶する四次曲線]\label{nft:cor:quartic-memory}
次数一三様体側で、二重積の線型同型類は $B_4$ の線型座標変換と
非零スカラー倍による同値類にちょうど対応する。零四次式も許す。
$V(B_4)$ が滑らかな場合、固定した二重積を持つ圏同値類の族の次元は $28$ であり、
$B_4=0$ では $19$ である。
\end{corollary}
\begin{proof}
二重積は全射なのでその核だけを考えればよい。定理\ref{nft:thm:triple-kernel}で
得た二重平面の反標準線束により $v\mapsto Mv$, $w\mapsto aw+Q_2(v)$ の
形に限られる。小行列式を除いた二次関係は重み四の直線
$\langle 3w^2+B_4\rangle$ である。二つの規格形の $w$ の一次項を比較すれば
$Q_2=0$、残りの係数を比較すれば $B'_4(Mv)=a^2B_4(v)$ となる。
この条件は逆に二重積の同型を与える。

滑らかな四次曲線を固定すると $B_6$ は $28$ 次元空間の非空開を動く。
実際、$v\ne0$ では六次式の自由係数は局所での値と平面二方向の一次 jet を
任意に取れる。$H_w=0$ はその上の有限二重被覆なので、$H=0$ と二方向の
偏微分消滅による三条件が同時に起きる係数は稠密でない。
$v=0$ の $Z_H$ 上では $s,w$ がともに非零で滑らかである。
三次重みの座標の変更は平方完成した式の間では $s$ の再尺度化に限られる。
重み付き dilation を除いた $(M,a)$ の自由度は、種数三の滑らかな曲線の
有限自己同型群と $a$ の平方根の有限選択のみである。従って開集合の有限商の
次元は $28$ である。$B_4=0$ なら $H$ の孤立性は $B_6$ が滑らかな
平面六次曲線を定めることと同値であり、$w$ の尺度で $B_6$ の全体の
スカラーを吸収できる。よってその族は滑らかな平面六次曲線の射影同型類で
次元 $27-8=19$ である。これらの方程式の同型類と圏同値類の一致は
定理\ref{nft:thm:triple-kernel}による。
\end{proof}

\subsection{次数一三様体の被覆対合を持ち上げる形式変形}
\label{nft:sec:deck-formal}
\begin{theorem}\label{nft:thm:deck-formal}
任意の滑らかな $Z_H$ について、$\mathcal C=\Ku(Z_H)$ とし、
$\tau$ を $s\mapsto-s$ による自己同値とする。
特殊ファイバーを固定し、通常の局所 Artin $\C$ 代数上で
定理\ref{thm:nf-residual-formal}の compactly generated な圏の変形を取る。
その形式 hull の同型類関手は $35$ 次元の形式円板であり、
$\tau$ の固定点部分は滑らかな $34$ 次元の形式円板である。
実際の圏の変形が $\tau$ の持上げを持つことは、方程式 $s^2+H_A$ の相対
残余圏で実現されることと同値である。持上げは存在のみを要求する。
\end{theorem}
\begin{proof}
$\mathcal J=S_0/J$ と置く。まず内部 twist 零で
\cite[Theorem~5.39 の証明]{NFBFKKernels}を適用する。
射の次数を $j$、移動空間の余次元を $c$、その重みの和を $b$ とすると、
$j-c$ が偶数のとき固定多項式の Jacobian 環の重み $b+3(j-c)$ を取る。
恒等元では偶数 $j$ に $\mathcal J_{3j}$ を得る。位数六では $(c,b)=(5,8)$ で
固定座標がなく、$8+3(j-5)=0$ の整数解はない。
位数三では $(c,b)=(4,5)$ で固定式は $s^2$、やはり重み零にならない。
位数二では $(c,b)=(4,6)$、固定式は $w^3$ である。
重み $3j-6$ が $0,2$ のいずれかになる偶数 $j$ は $j=2$ だけである。
$\tau$ は恒等 sector の $\mathcal J$ に恒等に作用する。位数二の sector では
法方向 $v_1,v_2,v_3,s$ の行列式双対から符号 $-1$ が付く。従って
\[
  HH^1(\mathcal C)=HH^3(\mathcal C)=0,\qquad
  HH^2(\mathcal C)^+=\mathcal J_6,\qquad
  HH^2(\mathcal C)^-=\C .
\]
実際すべての奇数次数が消える。$C^\bullet(\mathcal C)[1]$ の極小
$L_\infty$ 模型は奇数次数のみなので、各高次括弧の次数 $2-j$ により
$j$ 項括弧は全て零である。従って通常の Artin 基底上の形式 hull は主張の
形式円板である。対合で平均した形式座標を使えば、固定点部分は反不変座標を
零とした形式因子である（定理\ref{ifk:thm:rotation-formal}の固定点部分の議論）。

次に $(S_0)_6$ 内に $J_6$ の線型補空間を取り、その基底 $G_a$ で
$s^2+H+\sum t_aG_a$ を動かす。この滑らかな基点の形式近傍で、相対例外対と
その射影は係数変更と両立し、各点に $s\mapsto-s$ が作用する。従ってこの
$34$ 次元の形式族から hull への自然な射は固定点部分を通る。
この射の微分の単射性を周期から確かめる。重み付き留数表示
\cite[Lemmas~5.4--5.5]{SignLRZVeronese}において、
$G\in \mathcal J_6$ に沿う Kodaira--Spencer 作用は
$H^{2,1}\to H^{1,2}$ 上で、非零定数倍を除き
$\mathcal J_4\xrightarrow{G}\mathcal J_{10}$ である。
$\mathcal J_4 \mathcal J_4$ は $\mathcal J_8$ を張る。これは重み一、二の変数の重み八の単項式が
重み四ずつに分けられるためである。$\mathcal J_6$ と $\mathcal J_8$ の socle 対は完全なので、
この作用が零なら $G=0$ である。

この作用は残余圏への制限後も検出される。実際、包含による
$HH_{\pm1}(\mathcal C)\to HH_{\pm1}(Z_H)$ は例外対の加法性で同型であり、
HKR で右辺はそれぞれ $H^{2,1},H^{1,2}$ である。幾何的変形の
Hochschild 類を充満部分圏に制限すると、相対残余操作の接類になる
（定理\ref{thm:nf-residual-formal}の接写像の議論）。
この制限と包含による鎖の写像は cap 積と両立する。
充満部分圏内の対象列では制限した cochain の挿入が元の挿入と同じだからである。
幾何的接類の cap 積は HKR の上記 Hodge 成分で Kodaira--Spencer 縮約を誘導する。
ここでは $H^{1,0},H^{3,2}$ とそれらの共役が零で、
同じ次数差の他の成分による補正はない。
従って圏の接類の消滅から上記の乗法の消滅が従う。
固定点部分の接次元も $34$ なので形式逆写像定理で射は同型となる。

最後に $\operatorname{Ind}\mathcal C$ は proper な compact 部分圏で生成される。
定理\ref{thm:nf-residual-formal}の証明の tameness による hull 比較
\cite[Theorems~5.1.9, 5.3.16, Proposition~5.3.21]{NFDAGX}は
実際の変形の同型類上単射である。$\tau$ を持ち上げられる変形の hull 点は
固定されるので、上の形式族の一つの点と一致し、実際の $A$ 線型圏として同値になる。
逆向きには相対方程式の被覆対合を移送すればよい。これは特殊ファイバーを
指定した Artin 変形についての結論で、反不変方向の一般ファイバーの代数化や
可換三次五様体による実現は含まない。
\end{proof}

\begin{corollary}[偶符号族の三重積の退化因子]\label{nft:cor:reduced-product-divisor}
偶符号被覆の滑らかな係数開で、$\operatorname{Sym}^3 R_1\to R_3$ の
階数 $34$ 以下の退化スキームは被約因子 $\lambda=0$ である。
この因子上の各点を通る、他の係数を固定した $\lambda$ 方向は、
定理\ref{nft:thm:deck-formal}の次数一三様体の形式因子に横断する。
\end{corollary}
\begin{proof}
零係数の各点の近傍で $p_1,p_2,p_3$ を独立な線型座標 $u_1,u_2,u_3$ として取れる
（系\ref{nft:cor:allsigns-five}の証明により線型独立）。
相対残余圏の対角と相対内部 shift の間の Hom 複体は smooth proper 性から
perfect で、係数変更と両立する。滑らかな係数開の全点で $R_1,R_3$ の次元は
$7,35$ と一定なので、次数零はベクトル束となりその形成は任意の係数変更と両立する。
実際、局所自由複体の隣接二微分の階数は各々下半連続で、和が一定のため各々一定である。
基底は被約なので局所的に可逆小行列を消去すると該当項は直和因子になる。

$P=x_1x_2x_3$ による乗法は $R_3$ の切断である。源の方程式の偏微分から
恒等 sector で $x_i^2=-g_i$、$2u_ix_i^2+\lambda P=0$、従って
\[
                  \lambda P=2u_i g_i\quad (1\le i\le3)
\]
を得る。この式は各ファイバーで成立するので被約基底上の切断として成立する。
右辺は $u_1,u_2,u_3,w$ の三次式として
三重積の像に属する。各零係数点で三重積の像の $34$ 列を選ぶと、
これらと $P$ は局所的に $R_3$ の基底となる。選んだ列を除去した行の成分が
最大サイズの小行列式のイデアルを生成する。その全成分は $\lambda=0$ 上で
消え、この超平面が被約なので $(\lambda)$ に入る。逆に上式からこのイデアルは
$\lambda$ を含む。従って余核は局所的に $\mathcal O/(\lambda)$ であり、
非零係数では全射である。

次数一三様体の方程式を動かす滑らかな多項式族では全点で階数 $34$ なので、
同じ最大サイズの小行列式は恒等的に零である。そのどの Artin 底変更でも零のままである。
一方、偶符号族で $\lambda=\varepsilon$, $\varepsilon^2=0$ とするとこのイデアルは
$(\varepsilon)\ne0$ である。相対 Serre 関手からの積は相対同値で運ばれる
（基底線束の自明化の違いは小行列式のイデアルを変えない）。
よってこの一次変形は幾何的固定因子の接空間に入らず、その法線は一次元なので横断する。
\end{proof}
